\documentclass[conference]{IEEEtran}
\usepackage{cite}
\usepackage{amsmath,amssymb,amsfonts}
\usepackage{algorithmic}
\usepackage{graphicx}
\usepackage{textcomp}
\usepackage{xcolor}

\begin{document}

\title{Intelligent RF-VLC Hybrid Communication Framework for Autonomous Vehicles using Machine Learning \& Utility Optimization}

\author{\IEEEauthorblockN{Tejas Kondhalkar}
\IEEEauthorblockA{\textit{Electronics and Communication Engineering} \\
\textit{Faculty of Technology, University of Delhi}\\
New Delhi, India \\
tejasvilas04@gmail.com}
}

\maketitle

\begin{abstract}
This paper presents an intelligent hybrid communication framework that dynamically selects between Radio Frequency (RF) and Visible Light Communication (VLC) links for Vehicle-to-Vehicle (V2V) networks...
\end{abstract}

\section{Introduction}
Modern autonomous vehicles require ultra-reliable low-latency communication...

\section{System Model}
The proposed system utilizes physics-based channel modeling combined with machine learning...

% Formula Section
\subsection{RF Path Loss}
$PL = PL_0 + 10n\log_{10}(d) + shadow$

\section{Methodology}
We employ Random Forest and XGBoost classifiers to predict the optimal link...

\section{Conclusion}
The results demonstrate a 112.1\% utility improvement...

\end{document}
